\documentclass[fleqn,10pt]{olplainarticle}
% Use option lineno for line numbers 

\title{Pan-Cancer Multi-Agent Tumor Segmentation and Analysis with Conditional Probabilistic Anatomical Reasoning}

\author[1]{First Author}
\author[2]{Second Author}
\affil[1]{Address of first author}
\affil[2]{Address of second author}

\keywords{Keyword1, Keyword2, Keyword3}

\begin{abstract}
Please provide an abstract of no more than 300 words. Your abstract should explain the main contributions of your article, and should not contain any material that is not included in the main text. 
\end{abstract}

\begin{document}

\flushbottom
\maketitle
\thispagestyle{empty}

\section{Introduction}

\begin{itemize}
    \item Build a pan-cancer multi-agent tumor segmentation system
    \item 2D, 3D tumor segmentation models, anatomical segmentation models could be any, even allow users to define or use their private models
    \item Beyond segmentation, provide anatomical graph and explanation
\end{itemize}

\section{Methodology}
% Rewritten 16 Oct 2026 (Claude, with the user): ledger v6 (set-based host and tumor reasoning,
% documents/plan/Ledger_v6_rules_final_draft.md) and image grounding G0 v2 (documents/plan/Image_grounding_rules_draft.md).
% Previous version: documents/bak/main_before_g0_20261016.tex.
% Revised 18 Oct 2026: uncertainty-preserving host/tumor reasoning and one-pass conditional scoring.
% Pre-revision working manuscript preserved in documents/bak/main_before_reasoning_revision_20261018.tex.
\begin{figure*}[t]
    \centering
    \includegraphics[width=\linewidth]{figs/pipeline.pdf}
    \caption{Pan-cancer multi-agent tumor segmentation and reasoning. Given a scan $X$ and a user-given cancer type $C$,
    pluggable segmentation agents (defaults: BiomedParse for 2D tumor, VoxTell for 3D promptable tumor and anatomy,
    TotalSegmentator for anatomy; any model can fill a slot) produce tumor and anatomy masks. A planner turns the
    pan-cancer knowledge base $K$ into a scan-specific prompt plan $K_C$. Set-based reasoning turns the agents' masks into
    claims, anchors and alternative extents, preserving uncertainty in anatomical hosts. Frozen-reference image evidence
    and anatomical/model-support features enter a conditional score over existence, host and extent. Decision policies
    and lesion association are applied separately, with unresolved candidates retained for review.
    [legacy figure; to be redrawn for the revised design]}
    \label{fig:pipeline}
\end{figure*}

Figure~\ref{fig:pipeline} summarizes the pipeline; the components are detailed below.

\subsection{Conditional Probabilistic Anatomical Reasoning}
\label{sec:conditional_reasoning}
\begin{figure}
    \centering
    \IfFileExists{figs/PGM.png}{\includegraphics[width=1\linewidth]{figs/PGM.png}}{\fbox{\parbox{0.9\linewidth}{\centering figs/PGM.png not found}}}
    \caption{Anatomical reasoning graph. [to be redrawn: fixed uncertain anatomy $A$, frozen normal references, candidate evidence and
    conditional scores over the null state and joint tumor--host--extent alternatives; decision policy shown separately.]}
    \label{fig:pgm}
\end{figure}

We formulate tumor segmentation as conditional probabilistic anatomical reasoning that jointly infers each tumor, its
host organ and its extent, by combining the outputs of several segmentation agents with anatomical knowledge, set-based
logic and image evidence. Let $X$ be the image, $C$ the cancer type and $K$ the data-agnostic anatomical knowledge base.

\paragraph{Segmentation agents.}
Agents are pluggable and treated equally. A \emph{tumor agent} $k\in\mathcal{K}$ returns tumor masks for text prompts:
one cancer-specific prompt (for example \emph{kidney tumor} for KIRC) and, if the agent is promptable, one prompt per host
of the cancer type in the field of view. An \emph{organ agent} returns organ masks. With the default agents,
\begin{equation}
\label{eq:agents}
\begin{aligned}
B_C&=f_{\mathrm{BP}}(X,C), & V_C&=f_{\mathrm{VT}}(X,C,K_C),\\
V_K&=f_{\mathrm{VT}}(X,K_C), & T&=f_{\mathrm{TS}}(X),
\end{aligned}
\end{equation}
where $B_C$ and $V_C$ are tumor masks (VT's cancer prompt and its host prompts from the plan $K_C$,
Section~\ref{sec:kb_planning}), and $V_K$ and $T$ are organ masks. We write $A_{k,o}$ for organ agent $k$'s mask of
organ $o$ ($A_{\mathrm{TS},o}=T_o$, $A_{\mathrm{VT},o}=V_{K,o}$). Agreement counts agents, not prompts: different
prompts of one agent are one rater.

\paragraph{Design status and scope.}
This revised specification separates candidate generation, evidence assessment and decision policy. It supersedes the
earlier ledger-v6/G0 hard host assignments and EM-style refinement; implementation and evaluation of the revision are
pending. Anatomical compatibility and agent agreement are evidence features, not calibrated priors. The baseline uses a
frozen normal reference and a conditional scoring distribution, not a generative likelihood or an EM algorithm.

\paragraph{Anatomy with uncertainty.}
For organ $o$, let $\mathcal{K}_o$ be the agents that return technically valid nonempty masks, and let $\mathcal{J}_o$
contain other organs not related to $o$ by permitted containment in $K$. Define its possible territory
$U_o$, multi-agent interior $I_o$, disputed territory $C_o$ and uncertain boundary $D_o$ by
\begin{equation}
\label{eq:organ_reference}
\begin{aligned}
U_o&=\bigcup_{k\in\mathcal{K}_o}A_{k,o},\\
C_o&=\bigcup_{o'\in\mathcal{J}_o}(U_o\cap U_{o'}),\\
I_o&=\left(\bigcap_{k\in\mathcal{K}_o}A_{k,o}\right)\setminus C_o,\\
D_o&=U_o\setminus I_o .
\end{aligned}
\end{equation}
With no contributing agents, $U_o$ is empty and $I_o$ is unavailable; with one, $U_o$ is a single-source territory and
$I_o$ is again unavailable rather than independently confirmed. In either case no confident interior is asserted and
$D_o=U_o$. Agreement is not a guarantee of correctness. Mask quality, provenance and anatomical violations
are retained separately. Tumor inference never changes these masks: anatomy is fixed but explicitly uncertain. Missing
anatomy is distinguished from anatomy confidently outside the field of view. A coarse knowledge-base region can support
a missing host hypothesis, but cannot supply normal-tissue intensities.

\paragraph{Candidate units.}
A \emph{claim} is a connected tumor component with its model and prompt provenance. A claim is retained unless technically
invalid, such as an unrecoverable coordinate mismatch. Small size, flat shape, prompt mismatch and organ overlap are
warnings, not hard deletions. An \emph{anchor} is a shared core or a single-source candidate; unsupported residual components
of a claim are also retained as candidate units, with links to their parent claim. Their component construction and minimum
resolvable size must be frozen before evaluation. An \emph{outline} is a proposed extent linked to an anchor. Shared
provenance is recorded so that overlapping units are not interpreted as independent observations.

The candidate hosts $\mathcal{H}_a$ include anatomically compatible organs, coarse regions for expected but unsegmented
organs, and an unresolved-host state $\varnothing$. Proximity generates hypotheses, never a forced origin assignment.
Hosts outside the cancer-specific spread list can remain candidates if proposed tissue is compatible with their anatomy.
The spread list controls prompting priority, not the definition of tumor.

\paragraph{Candidate extents.}
Each anchor has a finite family of alternatives,
\begin{equation}
\label{eq:extents}
\mathcal{X}_a=\{a,\ \omega\ (\omega\in\Omega_a),
U_a=\bigcup_{\omega\in\Omega_a}\omega,
U_a\setminus D_a^{\mathrm{suspect}}\},
\end{equation}
where $D_a^{\mathrm{suspect}}$ is tissue flagged as possible leakage into other organs. The untrimmed alternative is always
retained: a leak flag can reflect real extension. Empty alternatives are removed. Duplicate or near-duplicate extents
(Dice at least $0.95$) are clustered by a deterministic rule; each cluster receives equal base mass and its members share
that mass. Prompt multiplicity therefore does not directly supply extra votes. Sensitivity to clustering and redundant
proposals is evaluated. No new boundary is synthesized beyond these recorded set operations.

\paragraph{Conditional scoring formulation.}
For each anchor the state space consists of one null state $0$ (no tumor supported) and states $(1,h,\xi)$ with
$h\in\mathcal{H}_a$ and $\xi\in\mathcal{X}_a$. The same observed candidate collection, anatomy and image are available to
all states. Let $f_a(h,\xi)$ contain anatomical compatibility, model-support and proposal-quality features, and let
$\ell_a(h,\xi)$ be a bounded image evidence score computed against a frozen normal reference. Define
\begin{equation}
\label{eq:conditional_joint}
s_a(h,\xi)=b+\beta^\top f_a(h,\xi)+\eta\,\ell_a(h,\xi).
\end{equation}
With fixed normalized base masses $m_a(h,\xi)>0$, equal across host hypotheses and extent clusters before evidence,
\begin{equation}
\label{eq:case_posterior}
\begin{aligned}
Z_a&=1+\sum_{h,\xi}m_a(h,\xi)\exp s_a(h,\xi),\\
R_a(0)&=Z_a^{-1},&
R_a(1,h,\xi)&=Z_a^{-1}m_a(h,\xi)\exp s_a(h,\xi).
\end{aligned}
\end{equation}
This is a discriminative conditional scoring distribution. It is not obtained by cancelling host-dependent normal
likelihoods, and normalized scores are not claimed to be calibrated clinical probabilities. The intercept and coefficients
are fixed on independent development cases, or declared as provisional heuristic settings if labeled development data are
unavailable; the latter setting supports ranking and review, not probability claims. No TCGA stage or final test ratings
may determine coefficients, feature selection or thresholds. With labeled development candidates, regularized conditional
log loss can fit the scores; reference-compatible hosts/extents must be specified before fitting, and all variants of a
patient stay in one split. Calibration requires separate held-out candidate-level validation.

\paragraph{Evidence provenance and dependence.}
Model agreement, silence, prompt identity and anatomical agreement are separate features. Multiple prompts from one model
are one source, and different models are not assumed independent. A model's tumor proposal checked against its own organ
mask is marked as same-source evidence. Image features can overlap information used by the segmentation agents; independence
is not asserted. Anatomical overlap enters $f_a$ once, not again as an independent image likelihood factor. Development
and ablation tests assess the incremental value of correlated features.

\paragraph{Frozen-reference inference.}
Build each normal reference once from reliable organ tissue outside all retained tumor proposals and a safety margin.
Missing normal tissue makes the corresponding image features unavailable; rejected candidates are not recycled into normal
tissue during inference. Evaluate Eq.~\eqref{eq:case_posterior} once, without updating case-specific weights from their own
responsibilities. For a candidate accepted by the decision policy, select a compatible joint state,
\begin{equation}
(\hat h_a,\hat\xi_a)=\arg\max_{h,\xi}R_a(1,h,\xi).
\end{equation}
Report existence support $p_a=1-R_a(0)$ and host/extent marginals conditional on $z_a=1$, together with their calibration
status. A maximum across overlapping candidate occupancy scores is only an occupancy summary, not a union probability.
Candidates with common claims remain dependent; their scores are not multiplied to produce lesion confidence.

\paragraph{Decision policy and cancer attribution.}
Acceptance, rejection and review thresholds are selected on independent development data for an explicit operating point
and frozen before testing. Missing evidence or ambiguous host/extent support can trigger review. Knowing that a primary
cancer exists does not force any candidate to be accepted; the scan can report \emph{primary not detected}. Tumor existence,
host assignment, cancer attribution and index-lesion selection are separate outputs. Anatomical compatibility may suggest
a primary-site candidate, possible contiguous extension or a secondary-site candidate; it does not establish invasion,
metastatic origin or another primary malignancy. Unresolved and out-of-scope abnormalities remain visible in a separate
review output. Any clinical override is recorded separately from the unmodified scores and segmentation decisions.

\paragraph{Lesion reporting.}
Candidate identity is independent of host inference. Nearby or overlapping candidates have association links, but proximity
alone does not merge them into one biological lesion. Duplicate detections are resolved with a frozen spatial/provenance
rule; ambiguous associations are flagged, and extent choices are retained for audit. Lesion counts and index-lesion
measurements are reported only after this association step. Association does not pool confidence or force a shared host.
A lesion-level probabilistic model, if introduced later, requires an explicit shared latent variable and a separate
evaluation.

\subsection{Pan-cancer knowledge base and prompt planning}
\label{sec:kb_planning}
% Drafted 4 Oct 2026 from knowledge base v0.5.7 (knowledge_base/*.yaml) and the planner (pancia_kb/planner.py,
% adapter.py). No equation of Section~\ref{sec:conditional_reasoning} is changed; this section defines how $K$ is
% built and how the prompt plan $K_C$ of Eq.~\eqref{eq:agents} is produced for each scan.

\begin{figure}
    \centering
    \includegraphics[width=0.9\linewidth]{figs/kb_graph.pdf}
    \caption{Pan-cancer anatomical knowledge graph}
    \label{fig:kb_graph}
\end{figure}

This section describes how the anatomical knowledge base $K$ is built once for all cancer types, and how, for each scan,
it is turned into the final VT prompt plan $K_C$ of Eq.~\eqref{eq:agents}. $K$ is a curated, image-independent graph;
$K_C$ is a finite list of prompt jobs, and $V_K=f_{\mathrm{VT}}(X,K_C)$ is the stack of masks VT returns for these jobs.
No component of $K$ or of the planner uses patient sex, and no rule or constant is specific to an organ, a cancer type
or a dataset.

\paragraph{Anatomical entities.}
$K$ contains 204 anatomical entities. Each entity has a category (solid organ, hollow organ, gland, bone, vessel,
lymph-node station, muscle, fascia, fat, duct, space, nerve, landmark or region) and a laterality (left, right, midline
or bilateral). It also records how it can be observed: the TotalSegmentator classes that represent it for CT and for
MR, if any; a VT main phrase with up to two aliases, each graded as in-vocabulary, near-vocabulary or
out-of-vocabulary against the published VT training vocabulary; its craniocaudal span between vertebral anchors;
whether it is always present, variably present or absent after common surgery; a typical volume range; and, for CT, a
typical Hounsfield-unit range. Finally, each entity that can host a tumor carries a tumor phrase (for example
\emph{liver tumor}). Thirty-nine entities are \emph{anchors}: structures that TotalSegmentator segments reliably and
that therefore fix the coordinate frame of a scan (vertebrae, major vessels and large organs). Landmarks are tied to
vertebral levels, so that every spatial prior in $K$ is expressed relative to the patient's own skeleton rather than in
millimetres.

\paragraph{Anatomical relations.}
Entities are connected by 433 typed, directed relations. Ten relation kinds are defined, of which nine are used:
\emph{adjacent\_to} (176), \emph{has\_part} (97), \emph{drains\_lymph\_to} (62), \emph{lies\_in} (23),
\emph{supplied\_by} (21), \emph{invested\_by} (16), \emph{drained\_by} (16), \emph{landmark\_for} (13) and
\emph{has\_duct} (9); the tenth, \emph{wall\_of}, is reserved. Each relation kind has a fixed spatial gate and a default
priority. The gate states where the target may lie relative to the source: inside it (\emph{has\_part}), in a shell
around it (\emph{invested\_by}), in a neighbouring box or directional shell (\emph{adjacent\_to}), in a corridor along
a vessel or duct (\emph{supplied\_by}, \emph{drained\_by}, \emph{has\_duct}) or in a body region
(\emph{drains\_lymph\_to}); \emph{lies\_in} and \emph{landmark\_for} carry no gate. Gates are later used to check VT
outputs against structures that are already segmented, and priorities order prompts when the budget is exceeded. The relations encode the anatomy that staging systems refer to: direct extension
into a neighbouring organ, invasion of an investing fascia or fat, and lymphatic drainage to a regional station.

\paragraph{Cancer-type spread sets.}
For each cancer type $C$, $K$ stores a weighted host set $\mathcal{H}_C$, which supplies prompting priorities, not an exhaustive restriction on the hosts considered in
Section~\ref{sec:reasoning_details}. Only two items are specified per cancer type: the primary organ (planning weight $1.0$) and a list of
distant metastatic sites (planning weight $0.15$). Organs at risk of local invasion (planning weight $0.3$) are not curated; they are
derived automatically as the entities reached from the primary organ by an \emph{adjacent\_to} or \emph{invested\_by}
relation. Each cancer type also has a short ensemble of tumor phrases (for example \emph{lung cancer}, \emph{lung
tumor}, \emph{lung mass}), and generic templates produce phrases for any host without a curated one. The distant-site lists are user-defined inputs rather than an exhaustive account of metastatic spread. Here we
restrict each list to the most common metastatic sites reported for that cancer type that are also anchors of $K$, and
cancer types without a curated list default to liver, lung, adrenal gland and spine. Head-and-neck anchors (brain, oral cavity,
pharynx, larynx and parotid gland) are part of $K$ but are in no default list, because the body scans considered here
rarely include the head and neck. Sites that are not anchors, such
as the pleura and distant lymph nodes, are not prompted. A user can extend a list, or add an anchor, without changing any other part of the framework, and only the affected cancer types need to be re-planned. The current
base covers 17 TCGA cancer types, plus a CPTAC entry whose cancer type is resolved per case. Adding a new cancer type
therefore requires only its primary organ, its metastatic sites and, optionally, its phrase ensemble; all other
knowledge is shared.

\paragraph{Planning $K_C$ for a scan.}
For a given scan, the planner builds $K_C$ deterministically in a single pass, without calling any segmentation model.
Its inputs are the modality, the cancer type $C$, a summary of the TotalSegmentator output $T$ (which classes are
present, their volumes and their extents), and the initial tumor proposals $B_C$ and $V_C$ of Eq.~\eqref{eq:agents}.
It proceeds in five steps.

\emph{(i) Frame.} The vertebrae found in $T$ define a craniocaudal ruler. The landmarks of $K$ are placed on this ruler,
and every entity span is converted to a slice range of the scan. This step also determines which entities lie inside the
field of view.

\emph{(ii) Anchors.} Each anchor is labelled \emph{found} (present in $T$ with a plausible volume and position),
\emph{expected} (inside the field of view but missing from $T$) or \emph{missing} (outside the field of view). Found
anchors give the reference structures against which the gates of later prompts are evaluated.

\emph{(iii) Tumor evidence.} A component of $B_C$ or $V_C$ is \emph{jointly proposed} when the two models agree on it; agreement is not validation. For each
host $h\in\mathcal{H}_C$ the planner records an evidence level $e_h\in\{0,1,2\}$: $2$ if a jointly proposed component lies in
or within a reach of $10$\,mm of $h$, $1$ if a component from one model does, and $0$ otherwise. It also records the
side of the primary tumor for paired organs. The planning weight of the host is $w_h(1+e_h)$, where $w_h$ is its
spread-set planning weight. The initial tumor proposals therefore influence only the order and the extent of the anatomical
prompts; they never remove a host from $\mathcal{H}_C$.

\emph{(iv) Prompt tiers.} Prompts are generated in four tiers. Tier~0 adds ruler prompts (spine, sacrum, hip bone) only
when $T$ provides neither vertebrae nor fallback landmarks, to attempt frame recovery; unresolved frame failures retain uncertain coverage. Tier~1 completes the anchors that are expected but missing from $T$, each gated by its
spatial prior. Tier~2 adds \emph{anatomical profiles} by walking the relation graph: two hops from the primary organ,
one hop from each secondary host, and the structures that can be reached from any host with tumor evidence
$e_h>0$. Tier~T contains the tumor prompts: one prompt with the tumor phrase of each host of $\mathcal{H}_C$ inside the
field of view. Tier~T is issued for every such host regardless of the evidence $e_h$, so that a lesion missed by both
initial proposals can still be proposed. For the same reason, when the initial VT tumor mask is empty or smaller than
$0.5$\,ml, the cancer-type phrase ensemble is added for the primary organ. This is a wording ensemble for the same
model, not a second rater: the initial and prompted VT outputs are counted as one model in tumor reasoning.

\emph{(v) Budget.} The plan is limited to $80$ prompts. Tiers~0, 1 and T are protected; only tier~2 is trimmed, lowest relation
priority first, with items derived from a host given a one-step priority boost.

The output $K_C$ is the list of prompt jobs that survive the budget. Each job records the entity, its side, the phrases to
issue (the main phrase and its first two aliases), its tier, its gate, and, for tumor prompts, the host and the host
planning weight $w_h$. VT is run once per scan on the full volume with all jobs of $K_C$; there is no re-planning after inference.
For each job, available phrase masks are aggregated by a prespecified voting rule and the individual masks and prompt
provenance are retained. Multiple phrases remain one model source. Spatial gates produce compatibility flags, not hard
output deletion: a component with less than the legacy $80\%$ within its expected gate is retained as uncertain. Empty
outputs mean not detected, not confirmed absence. Variable or surgically removable anatomy requires contextual evidence
before absence is asserted. The organ outputs form $V_K$; tier-T outputs join the tumor proposals, tagged by prompted host.
These uncertainty-preserving rules supersede the legacy planner's output filtering and require implementation migration.

Two properties of this construction matter for the reasoning that follows. First, although $K$ is independent of the
image, $K_C$ is not: it depends on $T$ through the frame and the anchors, and on the initial tumor proposals through the
evidence levels. This dependence is restricted to which structures are prompted and in what order, so it changes the
coverage of $V_K$ but not the meaning of any mask. Second, because tier~T prompts every plausible host whether or not it
contains a tumor, VT is regularly asked to segment a tumor that is absent. Its output for such a prompt can be a false
positive inside the prompted host or elsewhere in the scan. These outputs are therefore proposals rather than evidence that a tumor exists, and
their acceptance follows joint evidence scoring and the separate decision policy
(Section~\ref{sec:grounding_details}).

\subsection{Reasoning details}
\label{sec:reasoning_details}
% Revised 18 Oct 2026: uncertainty-preserving candidate reasoning and frozen-reference conditional scoring.
% This is a design specification; ledger_v6s.py and grounding_g0.py have not been migrated by this manuscript edit.

All geometry is measured in millimetres with verified image/mask alignment. Acquisition orientation and voxel spacing are
recorded. The following procedure is a proposed replacement for the earlier rule-led ledger, not a claim that the existing
implementation already follows it. Numerical feature thresholds inherited from G0 are provisional until independently
validated for this revised candidate population.

\paragraph{Order of computation.}
(0) image and coordinate quality checks; (1) organ uncertainty and field-of-view assessment; (2) claims, shared cores,
residual components and outline alternatives; (3) candidate hosts and evidence features; (4) frozen normal references and
image features; (5) one-pass conditional scoring; (6) acceptance/review policy, duplicate association and lesion reporting.
No earlier candidate is accepted merely to provide a reference for judging later candidates.

\paragraph{Image quality control.}
\label{sec:qc_details}
Each scan is graded good, degraded or technically unusable. Missing anatomy is not itself evidence that the image is
unusable. Expected-organ segmentation failures, noise and slice spacing are recorded separately, because a dedicated
protocol or tumor replacement can explain missing masks. Degraded scans retain candidate reasoning, with unreliable image
features marked unavailable rather than interpreted as evidence against tumor. Hard exclusion is reserved for unusable
image data or unrecoverable coordinate errors. Diagnostic thresholds and their validation must be finalized before the
revised cohort run; previous G0 quality-control labels are not silently reused as validation of this design.

\paragraph{Host reasoning.}
\label{sec:host_details}
\emph{Mask quality.} Implausible volume, laterality, vertebral position, component count, containment or adjacency produces
a quality flag for every host, including both primary and secondary sites. The flags retain measured severity and source.
Normal-anatomy violations can reflect disease or segmentation error, so masks are not clipped to another model's uncertain
parent mask. Only technically invalid masks are removed from spatial calculations, with reasons recorded.

\emph{Field of view and missing anatomy.} Organs are labeled observed, expected-but-unsegmented, uncertain coverage, or
confidently outside the scan. Coverage uncertainty propagates when the vertebral frame is unreliable. An expected but
missing organ has a coarse spatial hypothesis from $K$ for any host role; it has no normal-intensity reference. An organ
confidently outside the scan cannot host a visible candidate, but the unresolved-host state remains available.

\emph{Host generation.} Candidate hosts include organs whose possible territory or coarse region lies within a generous
search neighborhood of the candidate. The earlier $10$\,mm envelope is a provisional retrieval radius, not proof of origin;
sensitivity to this radius and candidate-host recall must be measured. Both the core and outlines can propose hosts, with
the source of the proposal recorded. Touching the primary never sets a probability to one. The unresolved-host state
remains available for every candidate.

\emph{Host evidence.} Features include fractions of the core and extent within confident and uncertain organ territory,
boundary contact, distance, competing-organ overlap, exophytic geometry, mask quality and cancer-site compatibility.
\begin{equation}
\label{eq:host_prior}
f_a^{\mathrm{host}}(h,\xi)=
(\text{core overlap},\text{extent overlap},\text{boundary contact},
\text{distance},\text{uncertainty},\text{site compatibility}).
\end{equation}
These are graded compatibility features, not independent likelihoods or a normalized anatomical prior. Distances and
overlaps unavailable for a missing organ are explicitly marked missing, not assigned values that imply incompatibility.
Knowledge-base role weights influence prompt priority; their values are not treated as calibrated probabilities of origin.

\paragraph{Tumor reasoning.}
\label{sec:tumor_details}
\emph{Claims and residual support.} Connected components retain model, prompt and host-prompt provenance. Shared cores and
single-source residual components are represented separately so that a correct shared core does not automatically validate
a large unsupported extension. Small components below a prespecified resolvability threshold are retained in the audit
output as unresolved, not labeled normal. Fragmentation and duplicate associations are assessed before lesion counting.

\emph{Evidence, not hard gates.} Flatness, small size, organ resemblance, apparent leakage and host-prompt mismatch enter
the feature vector with anatomy confidence. A prompt is evidence about the model query, not confirmation of its answer.
Large invasive or exophytic alternatives remain in the extent family despite these flags. Components are not rejected
solely for lying outside a curated metastatic-host list.

\emph{Model support.} Agreement is recorded by distinct model identity, with prompt multiplicity and same-model
tumor/anatomy provenance tracked separately. No fixed values of $0.9$, $0.75$ or $0.1$ are interpreted as existence
probabilities. Silence of another model means missing corroboration, not positive evidence for a particular candidate.
A model's incomplete sensitivity can be assessed on development data without forcing its silent output to be a negative
vote. The normal, tumor and indeterminate possibilities must remain distinguishable in review.

\emph{No primary-first admission.} Every candidate is scored before index selection. Multifocal candidates are not rejected
merely because a larger or better-supported primary was selected first. Adjacency to another candidate is an association
feature, not proof of invasion; host-specific agreement is not proof of metastatic origin. If no candidate clears the
acceptance policy, primary detection failure is reported even when a cancer diagnosis is known.

\paragraph{Image grounding.}
\label{sec:grounding_details}
\emph{Choose, not redraw.} Grounding ranks the retained extents. This bounds the spatial search but does not guarantee a
correct boundary: an agent can already have proposed normal tissue or missed part of a tumor. Extent-selection quality and
candidate-set oracle coverage are evaluated separately.

\emph{Frozen normal tissue.} Let $\hat H_h$ be the quality-screened multi-agent interior $I_h$, when available. A
single-source reference is an explicitly flagged fallback and is evaluated separately; a coarse region is never an
intensity reference. Let $P_{\mathrm{all}}$ be the union of all retained tumor proposals, before scoring. Define
\begin{equation}
\mathcal{N}_h=\hat H_h\setminus\operatorname{dilate}(P_{\mathrm{all}},2\,\mathrm{mm}).
\end{equation}
The reference is computed once and never expanded by rejected candidates. This conservative exclusion can leave too
little normal tissue; missing context is preferable to treating the candidate as its own normal reference. Within depth
bins ($0,1.5,3,5,8,12,18,25,35,50$\,mm and beyond), median and robust scale of the image smoothed with a $1$\,mm Gaussian
(in-plane for thick slices) define standardized deviation $Z(v)$. Depth is measured from the fixed organ reference, not
from a candidate-dependent organ union. Scale floors and handling of sparse bins must be fixed before evaluation.
A local reference requires at least $200$ voxels within $50$\,mm of the extent, with physical reference volume also reported;
otherwise host-frame signs are unavailable. These thresholds are provisional legacy settings.

\emph{Signs.} For an extent $\xi$ under host $h$, each sign is yes, no or unknown:
\begin{itemize}
\item \emph{Router} $S_0$: $\zeta=\operatorname{median}_{v\in\xi}Z(v)$; the extent deviates if $|\zeta|\ge1.5$.
\item \emph{Context} $S_1$ (closedness): the ratio of the volume of $\xi$'s same-sign deviation ($\operatorname{sign}(\zeta)Z>1.5$)
that stays connected to $\xi$ within $30$\,mm, inside $\hat H_h$ dilated by $2$\,mm, to the volume of $\xi$. Closed if
$\le1.42$, open if $>5.27$. Applies when $\xi$ deviates and at least half of $\xi$ lies inside $\hat H_h$ dilated by $2$\,mm.
\item \emph{Internal structure} $S_2$: the median local standard deviation ($2$\,mm) of the image standardised by $\xi$'s own
median and MAD, over $\xi$'s core. Patchy if $\le1.047$. Applies to extents of at least $3$\,ml.
\item \emph{Elongation} $S_3$: $\sqrt{\lambda_{\min}/\lambda_{\max}}$ of the voxel coordinates. Elongated if $<0.363$.
\end{itemize}
These numerical thresholds are inherited provisional G0 settings: the previous draft derived closedness and patchiness
from normal structures and open/elongated thresholds from expert tumors. They must be refitted using development data
under the revised reference construction. $S_0$ and $S_1$ depend on the host; $S_2$ and $S_3$ do not. Low within-extent
intensity scale makes $S_2$ unavailable unless a validated scale floor is used. Anatomy-overlap features are handled
separately in the conditional score. Contrast routes the context test; it is not assumed independent of proposal generation.
No claim is made that these features encode information unavailable to the segmentation models.

\emph{Evidence state.}
The provisional G0 hierarchy combines appearance and geometry:
\begin{equation}
\label{eq:tree}
\begin{aligned}
\text{tumor-like}&\iff(\text{closed}\vee\text{patchy})\wedge\neg\text{open}\wedge\neg\text{elongated},\\
\text{organ-like}&\iff(\text{open}\vee\text{elongated})\wedge\neg\text{closed}\wedge\neg\text{patchy}.
\end{aligned}
\end{equation}
The outcome is uncertain otherwise. Unknown signs are not treated as false: an outcome is assigned only when the observed
signs determine it for all completions, otherwise it is uncertain. Elongation alone without evaluable context is insufficient
for an organ-like assignment. Anatomy overlap is excluded from this image evidence state; it is represented once in
$f_a^{\mathrm{host}}$ or the proposal-quality features. Thresholds and the hierarchy must be revalidated with the frozen
reference, rather than assuming that results from the earlier adaptive-reference scheme transfer.

\emph{External evidence tables.}
\begin{equation}
\label{eq:likelihoods}
\ell_a(h,\xi)=\log\frac{\tilde p(E_a(h,\xi)\mid T)}
{\tilde p(E_a(h,\xi)\mid N,\text{organ class},\text{modality},\text{size})}.
\end{equation}
Tumor reference outcomes come from expert tumor masks and prespecified outline perturbations; normal reference outcomes
come from expert normal structures that resemble candidate proposals. Both classes use identical preprocessing,
reference construction and stratification. A log-ratio is a feature of the conditional score, not a Bayes factor asserted
to be independent of candidate generation. Tumor and normal tables are matched by organ class, modality and extent size,
with prespecified pooling for sparse strata, smoothing and bounded log-ratios. The legacy $[1/100,100]$ ratio bounds are
provisional sensitivity settings, not guarantees. Zero counts in both classes require explicit positive smoothing or a
neutral score, not a ratio of zeros. Fully unavailable appearance evidence contributes zero log-score and a missingness
flag; partial evidence uses a development-validated reduced table, or remains unavailable.

All variants of an expert mask remain in the same patient-level fold, and uncertainty is estimated by patient rather than
by perturbation count. External normal structures are not assumed representative of model-selected false positives:
candidate-level validation on independent agent outputs is required before probability or operating-point claims.
Calibration, preprocessing and threshold fitting use development partitions only. Through-plane harmonization, if used,
must record orientation and use a block factor of at least one; its effect on small lesions is evaluated.

\emph{Inference and decision.}
The frozen-reference baseline evaluates Eq.~\eqref{eq:case_posterior} once. Host and extent are selected jointly from
$R_a(1,h,\xi)$, not from separately re-estimated weights. No lesion-volume pooling or posterior-to-prior feedback is used.
Acceptance and rejection use separately frozen thresholds, with an intervening review region and explicit flags for
unavailable evidence or ambiguous host/extent assignments. A degraded scan is not silently subjected to a threshold
validated only for fully observed image features: its operating point must be validated separately or it remains review-only.
No candidate is immune to rejection merely because it is the largest primary-site proposal. If a user applies a clinical
override, both the score-based result and the override are retained.

\emph{Optional future refinement.}
Adaptive normal-reference updating is a separate experimental extension, not the default. It must retain a stable normal
reference component, exclude each tested candidate from its own reference, and monitor changes in existence scores,
reference membership, evidence states and selected extents. Initialization sensitivity, cycling and termination at an
iteration cap must be reported. Such a procedure is called alternating refinement unless a fixed objective and justified
updates are derived; neither EM monotonicity nor calibrated uncertainty is claimed.

\emph{Output.}
Report accepted target-cancer-compatible candidates separately from indeterminate or out-of-scope abnormalities. For each
candidate retain the full host/extent alternatives, joint selected state, uncalibrated or calibrated score status, evidence
provenance, anatomical uncertainty and policy flags. A target-compatible label is not pathological confirmation of cancer
origin. Duplicate association occurs before lesion counts and index measurements, with unresolved associations flagged.
The final mask is the union of selected accepted extents; overlapping scores are not multiplied and are not interpreted as
independent evidence of existence.

\section{Experiments and Results}
\label{sec:results}
% Evaluation plan revised 18 Oct 2026. Existing cohort counts are retained, not newly verified.
% No revised-method results have been generated. Stage is descriptive only.

\subsection{Data}
We use the public TCGA radiology collections of The Cancer Imaging Archive (TCIA). The cohort contains 3{,}291
baseline-study series (1{,}407 CT, 1{,}884 MR) from 15 cancer types: BLCA, BRCA, CESC, COAD, ESCA, KICH, KIRC, KIRP,
LIHC, LUAD, LUSC, OV, PRAD, STAD and UCEC. TCIA states that only the baseline study precedes surgery, and post-surgical
studies cannot be identified, so all evaluations use the baseline study only. The stage-evaluation scope is restricted
to baseline series with the primary organ in the field of view: 3{,}116 series from 942 patients. Scans graded \emph{technically unusable} by image quality control (Section~\ref{sec:qc_details}) are removed and reported with their reason ([TBD] series); \emph{degraded} scans ([TBD]) are retained with missing image evidence and a separately validated decision policy, or for review. Clinical stage comes
from the TCGA clinical records. No voxel-level tumour annotation exists for this cohort.

External image evidence tables are to be rebuilt from public expert annotations: tumour masks of the R2Seg collections
(bladder, breast, cervix, colon, kidney, liver, lung, pancreas, prostate and uterus, including benign uterine fibroids;
[TBD] masks) and normal anatomy of
TotalSegmentator CT (v3.0, scans without pathology) and MR (v3.0, scans from the original release source; scans indexed from the
Imaging Data Commons, which hosts TCIA collections, are excluded to avoid overlap with the evaluation cohort), adults only ([TBD] scans). Neither enters the evaluation cohort.

\subsection{Implementation status and development protocol}
The tumor agents are BiomedParse (2D) and VoxTell (3D promptable); the anatomy agents are TotalSegmentator and VoxTell. The
knowledge base has 204 entities and 433 relations (v0.5.7). The revised candidate retention, uncertain-host representation,
frozen references and conditional scoring are a proposed specification; the legacy ledger/G0 implementation and previous
outputs do not establish its performance. Runtime and cohort results remain [TBD].

Before implementation is frozen, specify the residual-component and duplicate-association rules, resolvability threshold,
host retrieval radius, feature definitions and scaling, extent clustering, sparse-reference handling, score coefficients,
and acceptance/review/rejection thresholds. These choices are versioned and developed independently of the final evaluation
patients. Without independent labeled candidate data, use declared heuristic scores and a review output; do not claim
calibration or a validated automatic decision policy. Record patient/source splits for table development, score fitting,
calibration and final testing, keeping all scan series and outline variants of one patient together.

\subsection{Evaluation without voxel ground truth}
\label{sec:eval_protocol}
Without voxel annotations, we use stage-free consistency indicators, not substitutes for segmentation accuracy. Different
sequences may show different lesion boundaries, and consistently wrong masks can appear stable. All are computed on
the baseline study only, with coverage and sequence differences recorded:
\begin{itemize}
  \item \emph{spread}: the median $|\log(v/\tilde v)|$ of the index-lesion volume $v$ across the series of one study,
  where $\tilde v$ is the study median;
  \item \emph{cross}: the share of same-study series pairs whose sizes differ by at least one RECIST progression step;
  \item \emph{disagree}: the share of series whose index lesion is not the study's majority lesion;
  \item \emph{miss}: the share of baseline series with no measurable lesion (a known cancer does not guarantee a visible or measurable primary on each series).
\end{itemize}
In addition, [TBD] agent candidates were rated blind by a reader (tumour, not tumour, unsure); these ratings measure the
precision of image evidence against a candidate and the share of rated tumours it rejects. The patient-median Spearman
correlation with T and overall stage is reported for reference only. No parameter, rule or threshold is selected on stage
or on the final test ratings. Any ratings previously used to select features belong to development and cannot also serve
as an independent test. Review and non-detection rates are reported alongside accepted-lesion precision, so abstention
cannot hide failures.

\subsection{Validation on expert data and independent candidates}
\label{sec:results_expert}
The revised pipeline is evaluated on patient-disjoint data with expert labels. Report: (i) candidate and candidate-host
recall before scoring, including missing-organ and exophytic cases; (ii) external evidence-table transfer to real
agent-selected tumors and false positives; (iii) joint host/extent selection accuracy and Dice relative both to the expert
mask and to the best available candidate extent; (iv) existence discrimination, false rejection, sensitivity and precision
at the frozen operating point; and (v) review coverage and performance stratified by organ, modality, size and anatomy
quality. If probability claims are made, report reliability curves, Brier score and log loss on independent candidate
labels, rather than on simulated perturbations alone. Uncertainty intervals resample patients.

Host/extent ambiguity and unavailable normal references are reported, not discarded. Annotation protocols must separate
tumor existence from attribution to the known cancer. External reference-table experiments use collection holdouts and
patient-level folds; transformations of a mask never cross splits. All results for the revised method remain [TBD].

\subsection{Effect of image-grounded reasoning}
\label{sec:results_r2}
The candidate set and non-image feature definitions are identical with and without image grounding. Compare the full score
against an image-feature ablation with its operating point selected on development data; also report fixed-coefficient
score changes to separate direct evidence effects from refitting.
Table~\ref{tab:r2_by_type} reports, per cancer type, the lesions admitted by reasoning only and after grounding; those
rejected and recovered by image evidence; anchors whose host or extent changed; and the volume between the anchor and the
union of outlines that the chosen extents keep. [TBD: main pattern.] On the rated candidates, image evidence against a
candidate is correct in [TBD]\,\% and rejects [TBD]\,\% of rated tumours.

\begin{table}[t]
\centering\small
\caption{Image-grounded reasoning per cancer type (baseline series). \emph{Reasoning only}: set-based host and tumour
features without image evidence. \emph{Rejected}/\emph{recovered}: matched candidates whose acceptance changes when image
features are included. \emph{Host}/\emph{extent}: candidates whose joint selected state changes. Review transitions and
association changes are reported separately. All cells are pending revised-method evaluation.}
\label{tab:r2_by_type}
\setlength{\tabcolsep}{4pt}\footnotesize
\begin{tabular}{lrrrrrrr}
\hline
Cancer & Series & \shortstack{Reasoning\\only} & Grounded & Rejected & Recovered & \shortstack{Host\\changed} & \shortstack{Extent\\changed}\\
\hline
\multicolumn{8}{c}{[TBD]}\\
\hline
All & 3{,}291 & & & & & & \\
\hline
\end{tabular}
\end{table}

\subsection{Baseline yardsticks}
\label{sec:results_yardsticks}
Table~\ref{tab:yardsticks} compares BiomedParse alone, VoxTell alone, reasoning only, and the full model on the
yardsticks of Section~\ref{sec:eval_protocol}. [TBD: main finding, stated without reference to stage.]

\begin{table}[t]
\centering\small
\caption{Stage-free yardsticks on 3{,}116 baseline series (942 patients). The Spearman $\rho$ with T stage is reported
only and was not used for any choice.}
\label{tab:yardsticks}
\begin{tabular}{lrrrrr}
\hline
Method & Spread & Cross (\%) & Disagree (\%) & Miss (\%) & $\rho_T$ (reported)\\
\hline
BiomedParse & & & -- & & \\
VoxTell & & & -- & & \\
Reasoning only & & & & & \\
Full model & & & & & \\
\hline
\end{tabular}
\end{table}

\subsection{Qualitative results}
[TBD: Fig.~\ref{fig:grounding} to be regenerated for the revised design: candidate alternatives, anatomical uncertainty,
frozen normal tissue, evidence scores, joint host/extent choice and acceptance/review policy. Existing G0 images are legacy illustrations.]
\begin{figure*}[t]
    \centering
    \IfFileExists{figs/r2_grounding_g0.pdf}{\includegraphics[width=\linewidth]{figs/r2_grounding_g0.pdf}}{\fbox{\parbox{0.9\linewidth}{\centering figure to be regenerated}}}
    \caption{Image-grounded reasoning on a public TCGA scan. Left: anchors (filled), outlines of each agent and the reasoning
    features. Middle: frozen references and conditional scores for host/extent alternatives. Right: accepted tumors and
    separate unresolved candidates, with joint selected extents. [Legacy illustration; revised outputs TBD]}
    \label{fig:grounding}
\end{figure*}

\subsection{Ablation}
\label{sec:ablation}
The primary baseline is fixed-reference, one-pass inference. Planned comparisons are: (i) each feature group removed
(model agreement/provenance, host compatibility, context, internal structure, geometry); (ii) uncertainty-preserving
candidate retention versus legacy hard gates; (iii) joint host/extent selection versus separate marginal choices;
(iv) frozen references versus a separately specified adaptive-reference extension; and (v) explicit unresolved-host and
review states versus forced decisions. Legacy EM-style refinement is a historical comparator, not the proposed algorithm.
Measure sensitivity to retrieval radius, extent redundancy, candidate fragmentation, anatomy perturbations and coefficient
choices. For any iterative extension also report initialization sensitivity, cycles, reference contamination and iteration
cap termination. All comparisons include candidate recall, false rejection, extent quality, calibration where supported,
and review coverage. [TBD]

\subsection{Limitations}
\begin{itemize}
  \item The revised method is not yet implemented or validated. Conditional scores require independent candidate-level
  calibration before interpretation as probabilities; retained correlated proposals and model errors remain important risks.
  \item Conservative frozen references may exclude too much tissue, particularly when proposals cover most of an organ.
  Missing-context and review rates are part of the reported performance, not hidden exclusions.
  \item Image quality control relies on what the anatomy agents can segment; a scan whose organs fail for reasons other than quality (unusual anatomy, post-surgical change outside the primary, or a dedicated protocol with little anatomy in view) may be graded degraded.
  \item There is no voxel-level ground truth, so accuracy is assessed through stage-free yardsticks, expert data and
  blind ratings of candidates.
  \item Image evidence decides existence and chooses among the agents' outlines; it does not redraw boundaries. A tumour
  that no agent outlines, or outlines only in part, is not recovered beyond the agents' proposals.
  \item Many agent false positives remain uncertain rather than organ-like: normal sub-structures that resemble tumour
  (intra-organ vessels, renal sinus, bowel loops) are the main confounder, and iso-intense candidates have no context test.
  \item Evidence transfer across CT/MR, protocols and tumor sites requires separate validation. Hosts without expert normal
  anatomy (for example ovary and breast) require pooled tables or unavailable-evidence handling.
  \item The default distant-site lists do not include brain, and pleura and distant lymph nodes have no anchor in
  the knowledge base, so proposal coverage at these sites is limited. Unexpected proposals remain unresolved rather than
  being declared normal, but this does not recover lesions that were never proposed.
\end{itemize}

\section*{Acknowledgments}

Additional information can be given in the template, such as to not include funder information in the acknowledgments section.

\bibliography{modular_refs}

\end{document}
